\clearpage
\section{Wymagania funkcjonalne i niefunkcjonalne}		%2

W rozdziale przedstawiono wymagania dla systemu wykrywania intruzów opartego na analizie obrazu z kamer, plików wideo oraz strumieni sieciowych. Celem systemu jest wspomaganie ochrony obiektu poprzez automatyczne wykrywanie obecności osób, ocenę zgodności ubioru z przyjętym wzorcem stroju firmowego, archiwizację zdarzeń oraz szybkie przekazywanie informacji operatorowi.

Wymagania podzielono na funkcjonalne i niefunkcjonalne. Pierwsza grupa opisuje czynności, które aplikacja musi wykonywać z punktu widzenia operatora, na przykład obsługę kamer, klasyfikację osób i zapis zdarzeń. Druga grupa dotyczy jakości działania systemu: wydajności, stabilności, czytelności interfejsu, możliwości konfiguracji oraz dalszej rozbudowy. Taki podział ułatwia późniejszą ocenę, czy system spełnia zarówno założenia użytkowe, jak i techniczne.

\subsection{Wymagania funkcjonalne}
\begin{itemize}
\item System musi obsługiwać wiele źródeł obrazu, w szczególności kamery lokalne, pliki wideo oraz strumienie sieciowe.
\item System musi umożliwiać skanowanie dostępnych kamer oraz dodawanie źródeł obrazu z poziomu interfejsu graficznego.
\item System musi wykrywać osoby w obrazie z wykorzystaniem modelu detekcji obiektów oraz umożliwiać ich śledzenie pomiędzy kolejnymi klatkami.
\item System musi pracować w dwóch trybach: dziennym oraz nocnym, przy czym wybór trybu może być realizowany automatycznie na podstawie harmonogramu lub ręcznie przez operatora.
\item System musi wyznaczać dla wykrytej osoby obszar analizy ubioru na podstawie maski segmentacyjnej sylwetki, tak aby możliwe było określenie koloru górnej i dolnej części stroju.
\item System musi umożliwiać zdefiniowanie wzorca koloru stroju firmowego, do którego porównywany jest ubiór wykrytej osoby.
\item W trybie dziennym system musi oznaczać jako potencjalnego intruza osoby, których ubiór nie odpowiada przyjętemu wzorcowi ubioru firmowego.
\item W trybie nocnym system musi traktować każdą wykrytą osobę jako potencjalny incydent, niezależnie od koloru ubioru.
\item System musi umożliwiać definiowanie masek ignorowanych obszarów kadru dla wybranych źródeł obrazu.
\item System musi pozwalać na konfigurację progów detekcji, parametrów alarmu oraz zasad archiwizacji zdarzeń z poziomu interfejsu operatora.
\item System musi umożliwiać pobieranie i ładowanie różnych wariantów modeli YOLO z poziomu aplikacji, bez konieczności modyfikacji kodu źródłowego.
\item System musi umożliwiać zapisywanie i wczytywanie presetów ustawień.
\item System musi rejestrować informacje o incydencie, w szczególności czas wykrycia, identyfikator źródła, tryb pracy, liczbę wykrytych osób oraz zapisany materiał dowodowy.
\item System musi archiwizować zdarzenia w postaci klipu wideo, gdy osoba pozostaje widoczna przez zdefiniowany czas.
\item System musi automatycznie ograniczać liczbę zapisanych zdarzeń zgodnie z limitem określonym w konfiguracji.
\item System musi prezentować operatorowi podgląd na żywo z wielu źródeł oraz sygnalizować stan alarmowy dla wybranego źródła.
\item System musi umożliwiać przegląd historii wykrytych zdarzeń wraz z podglądem zapisanego materiału.
\item System musi umożliwiać odtwarzanie nagrań z plików wideo, przewijanie materiału oraz powiększanie obrazu.
\end{itemize}

\subsection{Założenia użytkowe i techniczne}

System został zaprojektowany z myślą o pracy na stanowisku operatora ochrony. Oznacza to, że najważniejsze operacje powinny być dostępne z poziomu interfejsu graficznego: dodanie kamery, zmiana modelu, ustawienie wzorca ubioru, przegląd historii zdarzeń oraz eksport logów. Operator nie powinien być zmuszony do ręcznej modyfikacji kodu źródłowego ani do uruchamiania oddzielnych skryptów tylko po to, aby zmienić podstawowe parametry pracy aplikacji.

Drugim założeniem była możliwość dopasowania obciążenia systemu do dostępnego sprzętu. Z tego powodu wymagania obejmują obsługę różnych wariantów modeli YOLO, konfigurację rozdzielczości inferencji, limitu FPS modelu, progów detekcji oraz liczby maksymalnych wykryć. Pozwala to porównywać warianty nastawione na dokładność z wariantami oszczędnymi, takimi jak model nano. W praktyce jest to istotne, ponieważ aplikacja może pracować zarówno na pojedynczym źródle obrazu, jak i na kilku kamerach jednocześnie.

Wymagania uwzględniają również trwałość danych zdarzeń. System powinien zapisywać tylko te informacje, które są potrzebne do późniejszej analizy incydentu: czas, źródło obrazu, tryb pracy, liczbę osób, liczbę intruzów oraz ścieżkę do materiału wideo. Jednocześnie aplikacja musi ograniczać liczbę zapisanych zdarzeń, aby katalog archiwum nie rozrastał się bez kontroli podczas długotrwałej pracy.

\subsection{Wymagania niefunkcjonalne}
\begin{itemize}
\item System powinien analizować obraz w czasie zbliżonym do rzeczywistego, tak aby opóźnienie pomiędzy pojawieniem się osoby a jej wykryciem było możliwie małe.
\item System powinien zapewniać wysoką skuteczność detekcji osób oraz ograniczać liczbę fałszywych alarmów, zwłaszcza w trybie dziennym opartym na analizie koloru ubioru.
\item System powinien być odporny na długotrwałą pracę ciągłą i zachowywać stabilność podczas jednoczesnej obsługi wielu źródeł obrazu.
\item System powinien być skalowalny, aby możliwe było dodawanie kolejnych kamer lub strumieni bez konieczności przebudowy całej architektury.
\item System powinien umożliwiać łatwą rekonfigurację parametrów modelu, progów alarmowych, trybu pracy oraz ustawień archiwizacji.
\item System powinien przechowywać zdarzenia i logi w sposób uporządkowany, tak aby możliwy był ich późniejszy przegląd, audyt oraz analiza.
\item System powinien zapewniać modularną budowę, umożliwiającą dalsze rozszerzenie o bardziej zaawansowane metody segmentacji, analizy kolorów oraz klasyfikacji ubioru.
\item System powinien zapewniać bezpieczeństwo danych oraz ograniczenie dostępu do zapisanych nagrań, logów i konfiguracji.
\item Interfejs operatora powinien być czytelny, prosty w obsłudze i dostosowany do szybkiej oceny sytuacji alarmowej.
\item System powinien umożliwiać łatwe porównanie wariantów modeli pod kątem dokładności i obciążenia GPU.
\end{itemize}
